% Front matter: abstract, keywords, RQ, hypotheses, at-a-glance. Written LAST per plan.

\begin{abstract}
Retail inventory management turns on two coupled decisions: how much customers
will want, and how much stock to hold in anticipation. Most forecasting
research addresses only the first. This study addresses both. Twelve
forecasting approaches, from the moving average to a global long short-term
memory (LSTM) neural network, are compared on two retail demand environments:
500 intermittent-demand series from the Walmart M5 data and all 500 series of
a dense store-item demand dataset. Every model is evaluated under identical
conditions: the same calendar window, a 28-day horizon over eight rolling
origins, the same error metrics, and the same downstream inventory policy, so
that forecasts are judged by the costs they generate as well as the errors
they contain.

On sparse M5 demand, the LSTM records the lowest forecast error (mean absolute
scaled error, MASE, of 1.316) and the lowest inventory cost (152.83), ranking
first under 25 of 27 tested inventory policies. On dense store demand, the
LSTM again records the lowest forecast error (MASE 0.978) but ranks fourth on
inventory cost (2,247.46), behind a moving average (2,084.49); policy
leadership there is fragmented across moving average, SES, naive, and DES
(9, 7, 6, and 5 policies of 27, respectively). The implication is that lower
forecast error does not necessarily produce lower inventory cost. Ninety-one
paired statistical comparisons with Holm correction distinguish substantive
differences from sampling variation, and a 27-policy sensitivity grid
distinguishes conclusions that are stable across policies from ones that
depend on them. All data, code, and validation gates are frozen and traceable.
The large-language-model phase of this research program was designed but not
executed and appears here only as future work.
\end{abstract}

\noindent\textbf{Keywords:} retail demand forecasting; inventory optimization;
intermittent demand; M5 competition; LSTM; Croston method; forecast evaluation;
MASE; order-up-to policy; sensitivity analysis.

\subsection*{Research question}
How does the effectiveness of AI-based inventory optimization change as we move
from traditional forecasting models toward neural-network approaches, when
effectiveness is measured at the level of inventory outcomes rather than
forecast errors alone?

\subsection*{Hypotheses (neutral framing)}
\begin{description}[leftmargin=!,labelwidth=2.2cm]
\item[H1.] Forecast accuracy differs significantly across model generations on
the same series and horizons.
\item[H2.] Forecast-accuracy differences translate into measurable
inventory-outcome differences under a common policy.
\item[H3.] Increasing model sophistication yields monotonically improving
inventory outcomes. (The claim under test; rejection is a publishable
finding, not a failure.)
\item[H4.] Observed patterns are consistent across the two retail environments.
\end{description}

\subsection*{At a glance}
\begin{center}
\begin{tabular}{p{3.2cm}p{5.4cm}p{5.4cm}}
\toprule
 & \textbf{M5 (sparse, intermittent)} & \textbf{Store (dense, smooth)} \\
\midrule
Lowest forecast error & LSTM (MASE 1.316) & LSTM (MASE 0.978) \\
Lowest inventory cost & LSTM (cost 152.83) & Moving average (cost 2,084.49) \\
Stability & \textbf{Stable:} LSTM first in 25/27 policies & \textbf{Policy-dependent:} lead shared 9 / 7 / 6 / 5 \\
Correspondence & Lowest error $\rightarrow$ lowest cost & Lowest error $\not\rightarrow$ lowest cost \\
\bottomrule
\end{tabular}
\end{center}
\begin{center}
\emph{Lower forecast error does not necessarily imply lower inventory cost.}
\end{center}
